\section*{Capítulo \ref{ch:b3:hydrogeology} --- Hidrogeologia}

\begin{solution}{exo:b3:hydrogeology:1}
Primeiro filtro: $z = 120 - 40 = \qty{80}{m}$, $h = 120 - 12.5 =
\qty{107.5}{m}$, carga de pressão $107.5 - 80 = \qty{27.5}{m}$. Segundo
filtro: $z = \qty{40}{m}$, $h = \qty{105.0}{m}$, carga de pressão
\qty{65.0}{m}. A carga é maior no filtro raso, portanto a água
flui para baixo, sob um gradiente $(107.5 - 105.0)/40 = \num{0.0625}$:
os piezômetros estão em uma área de recarga. A pressão sozinha teria
sugerido o contrário.
\end{solution}

\begin{solution}{exo:b3:hydrogeology:2}
$q = K\Delta h/L = 12\times0.30/0.50 = \qty{7.2}{m/d}$. A área é
$\pi\times0.05^2 = \qty{7.85e-3}{m^2}$, logo $Q = \qty{0.0565}{m^3/d}$, ou
$56.5/24 = \qty{2.36}{L/h}$. A \omterm{def:b3:hydrogeology:seepage}{velocidade de percolação} é $7.2/0.35 =
\qty{20.6}{m/d}$.
\end{solution}

\begin{solution}{exo:b3:hydrogeology:3}
$K = 2.5/86400 = \qty{2.894e-5}{m/s}$, e $k = K\mu/(\rho g) =
\num{2.894e-5}\times\num{1.306e-3}/(999.7\times9.81) =
\qty{3.85e-12}{m^2}$. A \qty{20}{\celsius}, $K$ é multiplicado por
$(998.2/\num{1.002e-3})/(999.7/\num{1.306e-3}) = \num{1.30}$: $K =
\qty{3.25}{m/d}$. A variação se deve quase inteiramente à viscosidade.
\end{solution}

\begin{solution}{exo:b3:hydrogeology:4}
$T = 2.5\times180 = \qty{450}{m^2/d}$. A expansão da água dá
$\num{9e-7}\times180 = \num{1.62e-4}$, cerca de \qty{80}{\percent} de $S =
\num{2e-4}$; a compressão do calcário fornece o restante.
\end{solution}

\begin{solution}{exo:b3:hydrogeology:5}
Pela \cref{prop:b3:hydrogeology:thiem}, $T = Q\ln(r_2/r_1)/(2\pi(s_1 -
s_2)) = 1500\times\ln10/(2\pi\times1.53) = \qty{359}{m^2/d}$, e $K =
359/40 = \qty{9.0}{m/d}$.
\end{solution}

\begin{solution}{exo:b3:hydrogeology:6}
$Q/4\pi T = 1200/(4\pi\times300) = \qty{0.318}{m}$. Após 1 dia, $u =
300^2\times\num{5e-4}/(4\times300\times1) = \num{0.0375}$ e $W = -0.5772
+ 3.2834 + 0.0375 - 0.0004 = \num{2.743}$: $s = \qty{0.873}{m}$. Após 10
dias, $u = \num{0.00375}$, $W = -0.5772 + 5.5860 + 0.0038 = \num{5.013}$:
$s = \qty{1.596}{m}$. Dez vezes mais tempo dá menos que o dobro do
rebaixamento.
\end{solution}

\begin{solution}{exo:b3:hydrogeology:7}
$u \le 0.05$ quando $t \ge r^2S/(4T\times0.05)$. Para \qty{60}{m}: $3600\times
\num{2e-4}/90 = \qty{0.008}{d}$, \qty{11.5}{min}. Para \qty{500}{m}:
\qty{0.556}{d}, \qty{800}{min}, mais de 13 horas. O tempo cresce com $r^2$:
em um teste de um ou dois dias, os poços de observação usados para a reta
devem estar a algumas centenas de metros do poço bombeado.
\end{solution}

\begin{solution}{exo:b3:hydrogeology:8}
No ponto médio, cada poço está a \qty{200}{m} de distância: $u = 200^2\times\num{2e-4}/(4
\times450\times10) = \num{4.44e-4}$, $W = \num{7.142}$, $s = 0.382\times
7.142 = \qty{2.728}{m}$ de cada um, \qty{5.456}{m} no total. Para
\qty{400}{m}, $u = \num{1.78e-3}$, $W = \num{5.757}$, $s =
\qty{2.199}{m}$. Cada poço rebaixa a água do outro em
\qty{2.2}{m}: sua bomba precisa elevar a água \qty{2.2}{m} a mais, o que
custa energia, e o rebaixamento restante antes que a captação da bomba fique a seco é
menor.
\end{solution}

\begin{solution}{exo:b3:hydrogeology:9}
Quatro linhas de fluxo entre a base da barragem e o fundo delimitam $N_f = 5$
canais de fluxo; nove equipotenciais dividem a diferença de \qty{10}{m} em $N_d = 10$
quedas. $q = KHN_f/N_d = 5\times10\times5/10 = \qty{25}{m^3/d}$ por metro
de barragem, e \qty{7500}{m^3/d} sob \qty{300}{m}. O fator de
forma calculado, 0.512 em vez de 0.5, daria \qty{25.6}{m^3/d} por metro.
\end{solution}

\begin{solution}{exo:b3:hydrogeology:10}
$k = na^2/8 = 0.30\times(\num{5e-5})^2/8 = \qty{9.4e-11}{m^2}$, e $K =
k\rho g/\mu = \num{9.375e-11}\times998.2\times9.81/\num{1.002e-3} =
\qty{9.2e-4}{m/s}$, \qty{79}{m/d}, dentro da faixa de areia limpa, de 0.3 a
\qty{150}{m/d}. O modelo desconsidera a tortuosidade dos caminhos, mais longos
do que os tubos retos (o que reduz $K$), as gargantas entre os poros,
mais estreitas que os poros (também reduz), os poros que terminam em fundo cego e não transportam
fluxo (também reduz), e a dispersão dos tamanhos de poro: como um tubo transporta um
fluxo proporcional a $a^4$, os caminhos mais largos transportam a maior parte da água e
a dispersão eleva $K$ acima do valor de tubos de tamanho médio.
\end{solution}

\begin{solution}{exo:b3:hydrogeology:11}
(a)~O logaritmo se anula quando $2.25\,Tt = r^2S$, isto é, em $r_0 =
\sqrt{2.25\,Tt/S}$. (b)~Confinado e livre:
\[
r_0 = \sqrt{\frac{2.25\times450\times30}{\num{2e-4}}} = \qty{12.3}{km},
\qquad r_0 = \sqrt{\frac{2.25\times450\times30}{0.2}} = \qty{390}{m},
\]
$\sqrt{1000} = 31.6$ vezes menor. (c)~A pressão
propaga-se com a difusividade $T/S$, enorme quando $S$ é pequeno: dentro
do teste, o cone atinge falhas, os afloramentos do aquífero e os rios,
e o rebaixamento se afasta da reta, o que indica ao
hidrogeólogo que um limite foi atingido.
\end{solution}

\begin{solution}{exo:b3:hydrogeology:12}
(a)~Vista do referencial que se move com a água, a frente se espalha com
desvio-padrão $\sigma_x = \sqrt{2Dt} = \sqrt{2\alpha_Lvt}$. A frente passa
por $x$ por volta de $t = x/v$, quando $\sigma_x = \sqrt{2\alpha_Lx}$; à velocidade
$v$, um espalhamento de $\sigma_x$ em distância corresponde a um de $\sigma_t = \sigma_x/v$
no tempo. (b)~Os tempos de \qty{16}{\percent} e de \qty{84}{\percent} estão a um
desvio-padrão de cada lado da média: $\sigma_t = (247.2 -
152.6)/2 = \num{47.3}\text{ dias}$. Então $\alpha_L = v^2\sigma_t^2/(2x) =
0.25\times47.3^2/200 = \qty{2.8}{m}$. (A curva foi calculada com
\qty{3}{m}: a aproximação normal a subestima ligeiramente.) (c)~$R =
1 + 1.9\times0.32/0.30 = \num{3.03}$: o soluto se desloca a um terço da
velocidade da água.
\end{solution}

\begin{solution}{pb:b3:hydrogeology:1}
\textbf{1.} Xisto $363.2 - 282.5 = \qty{80.7}{m}$; calcário $282.5 -
102.5 = \qty{180}{m}$.

\textbf{2.} $h = 363.2 - 45.0 = \qty{318.2}{m}$; carga de pressão no topo
$318.2 - 282.5 = \qty{35.7}{m}$. A água sobe \qty{35.7}{m} acima do
topo do calcário, contida pelo xisto: o aquífero é confinado.

\textbf{3.} De $h_2 - h_1 = 550a - 100b = -1.62$ e $h_3 - h_1 = 250a
- 600b = -1.44$: $a = \num{-2.715e-3}$ e $b = \num{1.269e-3}$ por
metro.

\textbf{4.} $|\nabla h| = \sqrt{a^2 + b^2} = \num{0.0030}$. A água
flui ao longo de $-\nabla h = (\num{2.715e-3}, \num{-1.269e-3})$, para
SL: azimute $\ang{180} - \arctan(2.715/1.269) = \ang{115}$.

\textbf{5.} Os pontos anteriores a cerca de \qty{10}{min} situam-se acima
da reta: ali $u > 0.05$ e a curva de Theis fica acima de sua aproximação
logarítmica ($W > -\gamma - \ln u$, como mostra a figura da função de poço).

\textbf{6.} A melhor reta pelos pontos a partir de \qty{15}{min} passa por
\qty{1.75}{m} em \qty{100}{min} e \qty{2.63}{m} em \qty{1000}{min}:
$\Delta s = \qty{0.87}{m}$ por ciclo logarítmico.

\textbf{7.} $T = 2.3\times2160/(4\pi\times0.872) = \qty{453}{m^2/d}$.

\textbf{8.} A reta cai \qty{0.87}{m} por ciclo a partir de \qty{1.75}{m}
em \qty{100}{min}: o zero é atingido $1.75/0.872 = 2.01$ ciclos antes,
$t_0 = \qty{0.98}{min} = \qty{6.83e-4}{d}$.

\textbf{9.} $S = 2.25\,Tt_0/r^2 = 2.25\times453\times\num{6.83e-4}/3600 =
\num{1.9e-4}$, dentro da faixa dos \omterm{def:b3:hydrogeology:aquifer-types}{aquíferos confinados}.

\textbf{10.} Em \qty{15}{min}, $u = 3600\times\num{1.9e-4}/(4\times453
\times0.0104) = \num{0.036}$, e menor depois.

\textbf{11.} $K = 453/180 = \qty{2.5}{m/d}$, um calcário fraturado.

\textbf{12.} $u = 3600\times\num{2e-4}/(4\times450\times2) =
\num{2.0e-4}$, $W = -0.5772 + 8.5172 + 0.0002 = \num{7.940}$, $s =
0.382\times7.940 = \qty{3.03}{m}$, como registrado.

\textbf{13.} $u = 10^6\times\num{2e-4}/(4\times450\times30) =
\num{0.0037}$, $W = \num{5.02}$, $s = \qty{1.92}{m}$.

\textbf{14.} $r_0 = \sqrt{2.25\times450\times30/(\num{2e-4})} =
\qty{12.3}{km}$, mais que a largura da região. O cone atingiu as bordas
do calcário, seus afloramentos e a falha, que pode opô-lo a uma rocha
menos permeável: o aquífero infinito de Theis não o descreve mais, e
só o registro dos rebaixamentos dirá se um limite alimenta o poço ou
interrompe o fluxo.

\textbf{15.} $2160\times30 = \qty{64800}{m^3}$. Vem do armazenamento:
a água liberada pela expansão da água e pela compressão
do calcário, em toda a área do cone. Integrando $S\,s$ sobre
o cone devolve $Qt$.

\textbf{16.} $u = 10^6\times0.18/(4\times450\times30) = \num{3.33}$, $W
= \num{0.0086}$, $s = \qty{3}{mm}$: a fazenda não perceberia. O
bombeamento retiraria água dos poros próximos ao poço, em vez da pressão
de todo o aquífero.

\textbf{17.} $q = 2.5\times0.0030 = \qty{0.0075}{m/d}$; $v = 0.0075/0.18
= \qty{0.042}{m/d}$, \qty{15}{m} por ano.

\textbf{18.} $200/0.0417 = \num{4800}\text{ dias}$, \num{13.1}~anos; o
poluente, \num{32.9}~anos.

\textbf{19.} Por um cilindro de raio $r$ a vazão $Q$ atravessa a
área $2\pi rb$, da qual os poros são a fração $n_e$: $\mathrm dr/\mathrm
dt = -Q/(2\pi rbn_e)$, e integrando $2\pi bn_er\,\mathrm dr = -Q\,
\mathrm dt$ de $r$ a 0 obtém-se $t = \pi n_ebr^2/Q$. Aqui $\pi\times0.18
\times180\times200^2/2160 = \num{1885}\text{ dias}$, \num{5.2}\text{ anos} para a
água e $2.5\times5.2 = \num{12.9}\text{ anos}$ para o poluente: o poço o traz
mais de duas vezes mais cedo.

\textbf{20.} $\alpha_L = \qty{20}{m}$ e $\sqrt{2\times20\times200} =
\qty{89}{m}$. A concentração atinge \qty{16}{\percent} de seu
valor final \qty{89}{m} à frente da frente média: esse nível chega ao
poço quando a frente média ainda está a \qty{89}{m} de distância, em $(111/200)^2\times12.9 = \num{4.0}\text{ anos}$, um terço
do tempo médio. Os primeiros traços importam para a água potável, não
a média.

\textbf{21.} A favor: o poço captura o poluente, de modo que ele não
avança mais a jusante, e a água pode ser tratada onde é bombeada.
Contra: o bombeamento atrai o poluente para um poço de abastecimento
mais cedo do que ele chegaria a qualquer lugar de outro modo, e a água
precisa ser monitorada e tratada por décadas.

\textbf{22.} A \omterm{def:b3:hydrogeology:storage}{transmissividade} do calcário Varden é $T \approx
\textbf{\qty[text-series-to-math=true]{450}{m^2/d}}$ (\qty{453}{m^2/d}
do ensaio).
\end{solution}
